% GENERATED FILE. Edit canonical lessons, then run npm run generate:books.
% canonical-source-hash: fnv1a64:50c5f837b62c3aec
% canonical-lessons: ML-C376-ira, ML-C376-jivan, ML-C376-mrtadeham, ML-C376-parikkelkkuka, ML-C376-munnimarikkuka

\chapter{Lives Lost and Saved}
\label{ch:ml-lives-lost-and-saved}
\hlchaptermodality{376}

\begin{chapteropening}
\textbf{By the end of this chapter:} \emph{I can say a victim, life, a dead body, to get hurt and to drown.}

Complete the last lesson of chapter 376: I can say a victim, life, a dead body, to get hurt and to drown.
\end{chapteropening}

\section[\texorpdfstring{\ml{ഇര} (ira)}{ഇര (ira)}]{\ml{ഇര} (ira) — a victim; prey}

\label{lesson:ML-C376-ira}



\begin{quote}
Before the new one: say the Malayalam for to escape, then the Malayalam for to get stuck.
\end{quote}

\subsection*{You'll want to know: \ml{ഇര}}
\textbf{\ml{ഇര}} — \emph{ira} — \textquotedblleft{}a victim; prey\textquotedblright{}.

\begin{quote}
\ml{തട്ടിപ്പിന്റെ} \ml{ഇരകൾ} \ml{പോലീസിൽ} \ml{പരാതി} \ml{കൊടുത്തു}.
\end{quote}

\begin{quote}
\emph{taṭṭippinṟe irakaḷ pōlīsil parāti koṭuttu.}
\end{quote}

\begin{quote}
The victims of the fraud complained to the police.
\end{quote}

Plural: \ml{ഇരകൾ}, \emph{irakaḷ}.

First the prey of an animal; in the news, a victim.

\begin{cousinweb}
\textbf{\ml{ഇര}} (\emph{ira}) is native Dravidian: DEDR 490 lists Malayalam \emph{ira}, \textquotedblleft{}food of birds, snakes\textquotedblright{}, beside Tamil \emph{irai}, \textquotedblleft{}prey\textquotedblright{}. The news uses it for a person something bad happened to.
\end{cousinweb}

\begin{grammarlens}[title={What you've built}]
One more word for telling what happened.
\end{grammarlens}

\subsection*{Your turn}
\begin{itemize}
\raggedright
  \item \emph{Say it:} \emph{ira}
  \item \emph{Say it:} \emph{ira}, once more
  \item \emph{Say it:} \emph{taṭṭippinṟe irakaḷ pōlīsil parāti koṭuttu.} — The victims of the fraud complained to the police.
  \item \emph{Recall:} say the Malayalam for to escape, then the Malayalam for to get stuck, then say \emph{ira} again
\end{itemize}

\begin{tcolorbox}[breakable,colback=teal!4,colframe=teal!35!black,title={Before you move on}]
What does \ml{ഇര} mean? (\textquotedblleft{}A victim; prey\textquotedblright{}.) Say it once more.
\end{tcolorbox}
\section[\texorpdfstring{\ml{ജീവൻ} (jīvan)}{ജീവൻ (jīvan)}]{\ml{ജീവൻ} (jīvan) — life; a life}

\label{lesson:ML-C376-jivan}



\begin{quote}
Before the new one: say the Malayalam for to get stuck, then the Malayalam for a victim.
\end{quote}

\subsection*{You'll want to know: \ml{ജീവൻ}}
\textbf{\ml{ജീവൻ}} — \emph{jīvan} — \textquotedblleft{}life; a life\textquotedblright{}.

\begin{quote}
\ml{നഴ്സ്} \ml{കുഞ്ഞിന്റെ} \ml{ജീവൻ} \ml{രക്ഷിച്ചു}.
\end{quote}

\begin{quote}
\emph{naḻsŭ kuññinṟe jīvan rakṣiccu.}
\end{quote}

\begin{quote}
The nurse saved the baby's life.
\end{quote}

Saving a life is \ml{ജീവൻ} \ml{രക്ഷിക്കുക}. \ml{ജീവിക്കുക}, to live, has the same root.

\begin{cousinweb}
\textbf{\ml{ജീവൻ}} (\emph{jīvan}) is borrowed from Sanskrit \emph{jīva}, \textquotedblleft{}living; a living being, life\textquotedblright{} (CDIAL 5239), the root of \ml{ജീവിക്കുക} (\emph{jīvikkuka}), \textquotedblleft{}to live\textquotedblright{}.
\end{cousinweb}

\begin{grammarlens}[title={What you've built}]
One more word for telling what happened.
\end{grammarlens}

\subsection*{Your turn}
\begin{itemize}
\raggedright
  \item \emph{Say it:} \emph{jīvan}
  \item \emph{Say it:} \emph{jīvan}, once more
  \item \emph{Say it:} \emph{naḻsŭ kuññinṟe jīvan rakṣiccu.} — The nurse saved the baby's life.
  \item \emph{Recall:} say the Malayalam for to get stuck, then the Malayalam for a victim, then say \emph{jīvan} again
\end{itemize}

\begin{tcolorbox}[breakable,colback=teal!4,colframe=teal!35!black,title={Before you move on}]
What does \ml{ജീവൻ} mean? (\textquotedblleft{}Life; a life\textquotedblright{}.) Say it once more.
\end{tcolorbox}
\section[\texorpdfstring{\ml{മൃതദേഹം} (mṛtadēhaṁ)}{മൃതദേഹം (mṛtadēhaṁ)}]{\ml{മൃതദേഹം} (mṛtadēhaṁ) — a dead body}

\label{lesson:ML-C376-mrtadeham}



\begin{quote}
Before the new one: say the Malayalam for a victim, then the Malayalam for life.
\end{quote}

\subsection*{You'll want to know: \ml{മൃതദേഹം}}
\textbf{\ml{മൃതദേഹം}} — \emph{mṛtadēhaṁ} — \textquotedblleft{}a dead body\textquotedblright{}.

\begin{quote}
\ml{നാട്ടുകാർ} \ml{പുഴയുടെ} \ml{തീരത്ത്} \ml{ഒരു} \ml{മൃതദേഹം} \ml{കണ്ടു}.
\end{quote}

\begin{quote}
\emph{nāṭṭukār puḻayuṭe tīrattŭ oru mṛtadēhaṁ kaṇṭu.}
\end{quote}

\begin{quote}
The locals saw a dead body on the river bank.
\end{quote}

The respectful news word for a body; the plain old word is \ml{ശവം}.

\begin{cousinweb}
\textbf{\ml{മൃതദേഹം}} (\emph{mṛtadēhaṁ}) is built from two Sanskrit words, \emph{mṛta}, \textquotedblleft{}dead\textquotedblright{}, and \emph{deha}, \textquotedblleft{}the body\textquotedblright{} (Monier-Williams).
\end{cousinweb}

\begin{grammarlens}[title={What you've built}]
One more word for telling what happened.
\end{grammarlens}

\subsection*{Your turn}
\begin{itemize}
\raggedright
  \item \emph{Say it:} \emph{mṛtadēhaṁ}
  \item \emph{Say it:} \emph{mṛtadēhaṁ}, once more
  \item \emph{Say it:} \emph{nāṭṭukār puḻayuṭe tīrattŭ oru mṛtadēhaṁ kaṇṭu.} — The locals saw a dead body on the river bank.
  \item \emph{Recall:} say the Malayalam for a victim, then the Malayalam for life, then say \emph{mṛtadēhaṁ} again
\end{itemize}

\begin{tcolorbox}[breakable,colback=teal!4,colframe=teal!35!black,title={Before you move on}]
What does \ml{മൃതദേഹം} mean? (\textquotedblleft{}A dead body\textquotedblright{}.) Say it once more.
\end{tcolorbox}
\section[\texorpdfstring{\ml{പരിക്കേൽക്കുക} (parikkēlkkuka)}{പരിക്കേൽക്കുക (parikkēlkkuka)}]{\ml{പരിക്കേൽക്കുക} (parikkēlkkuka) — to get hurt; to be injured}

\label{lesson:ML-C376-parikkelkkuka}



\begin{quote}
Before the new one: say the Malayalam for life, then the Malayalam for a dead body.
\end{quote}

\subsection*{You'll want to know: \ml{പരിക്കേൽക്കുക}}
\textbf{\ml{പരിക്കേൽക്കുക}} — \emph{parikkēlkkuka} — \textquotedblleft{}to get hurt; to be injured\textquotedblright{}.

\begin{quote}
\ml{അപകടത്തിൽ} \ml{രണ്ട്} \ml{പേർക്ക്} \ml{പരിക്കേറ്റു}.
\end{quote}

\begin{quote}
\emph{apakaṭattil raṇṭŭ pērkkŭ parikkēṟṟu.}
\end{quote}

\begin{quote}
Two people were injured in the accident.
\end{quote}

Past: \ml{പരിക്കേറ്റു}, \emph{parikkēṟṟu}.

Literally to receive an injury: the person hurt takes -\ml{ക്ക്} or -\ml{ന്}.

\begin{cousinweb}
\textbf{\ml{പരിക്കേൽക്കുക}} (\emph{parikkēlkkuka}) joins \ml{പരിക്ക്}, \textquotedblleft{}an injury\textquotedblright{}, whose origin is not settled, and the native verb \ml{ഏൽക്കുക}, \textquotedblleft{}to receive, to take in\textquotedblright{} (DEDR 905).
\end{cousinweb}

\begin{grammarlens}[title={What you've built}]
One more word for telling what happened.
\end{grammarlens}

\subsection*{Your turn}
\begin{itemize}
\raggedright
  \item \emph{Say it:} \emph{parikkēlkkuka}
  \item \emph{Say it:} \emph{parikkēlkkuka}, once more
  \item \emph{Say it:} \emph{apakaṭattil raṇṭŭ pērkkŭ parikkēṟṟu.} — Two people were injured in the accident.
  \item \emph{Recall:} say the Malayalam for life, then the Malayalam for a dead body, then say \emph{parikkēlkkuka} again
\end{itemize}

\begin{tcolorbox}[breakable,colback=teal!4,colframe=teal!35!black,title={Before you move on}]
What does \ml{പരിക്കേൽക്കുക} mean? (\textquotedblleft{}To get hurt; to be injured\textquotedblright{}.) Say it once more.
\end{tcolorbox}
\section[\texorpdfstring{\ml{മുങ്ങിമരിക്കുക} (muṅṅimarikkuka)}{മുങ്ങിമരിക്കുക (muṅṅimarikkuka)}]{\ml{മുങ്ങിമരിക്കുക} (muṅṅimarikkuka) — to drown}

\label{lesson:ML-C376-munnimarikkuka}



\begin{quote}
Before the new one: say the Malayalam for a dead body, then the Malayalam for to get hurt.
\end{quote}

\subsection*{You'll want to know: \ml{മുങ്ങിമരിക്കുക}}
\textbf{\ml{മുങ്ങിമരിക്കുക}} — \emph{muṅṅimarikkuka} — \textquotedblleft{}to drown\textquotedblright{}.

\begin{quote}
\ml{ഭാഗ്യത്തിന്} \ml{ആരും} \ml{മുങ്ങിമരിച്ചില്ല}.
\end{quote}

\begin{quote}
\emph{bhāgyattinŭ āruṁ muṅṅimariccilla.}
\end{quote}

\begin{quote}
Luckily nobody drowned.
\end{quote}

Past: \ml{മുങ്ങിമരിച്ചു}, \emph{muṅṅimariccu}.

\ml{മുങ്ങുക}, to sink, + \ml{മരിക്കുക}, to die.

\begin{cousinweb}
\textbf{\ml{മുങ്ങിമരിക്കുക}} (\emph{muṅṅimarikkuka}) joins \ml{മുങ്ങുക} (\emph{muṅṅuka}), \textquotedblleft{}to sink, to go under\textquotedblright{}, and \ml{മരിക്കുക} (\emph{marikkuka}), \textquotedblleft{}to die\textquotedblright{}, which is built from Sanskrit \emph{maraṇa}, \textquotedblleft{}dying\textquotedblright{} (CDIAL 9869).
\end{cousinweb}

\begin{grammarlens}[title={What you've built}]
One more word for telling what happened.
\end{grammarlens}

\subsection*{Your turn}
\begin{itemize}
\raggedright
  \item \emph{Say it:} \emph{muṅṅimarikkuka}
  \item \emph{Say it:} \emph{muṅṅimarikkuka}, once more
  \item \emph{Say it:} \emph{bhāgyattinŭ āruṁ muṅṅimariccilla.} — Luckily nobody drowned.
  \item \emph{Recall:} say the Malayalam for a dead body, then the Malayalam for to get hurt, then say \emph{muṅṅimarikkuka} again
\end{itemize}

\begin{tcolorbox}[breakable,colback=teal!4,colframe=teal!35!black,title={Before you move on}]
What does \ml{മുങ്ങിമരിക്കുക} mean? (\textquotedblleft{}To drown\textquotedblright{}.) Say it once more.
\end{tcolorbox}
